%! TEX root = ../main.tex
\documentclass[main]{subfiles}

% \renewcommand{\appendixname}{付録} %%ここを付録と変更することで「付録 A」などと変更されます．デフォルトでは「appendix A」．
% \renewcommand{\appendixtocname}{付録} %%目次のappendix開始ページに表示される文字列を変更できます．
% \renewcommand{\appendixpagename}{付録} %%\appendixpage で表示されるページにでる文字列を変更できます．
% \appendixpage %%付録開始時に \appendixpagename と書かれたページを1枚追加します．
% \addappheadtotoc %%目次に付録が始まるページを追加してくれます．

% \begin{document}
% \chapter{付録}
% \section{概要}
% ソースコードは主に訓練から評価まで行うtrain\_pretrained.pyと，モデルのあらゆる定義を行うsed\_trainer\_pretrained.pyから成る．
% 付録では，CMTとMixStyleの主要な実装を示す．なお，SEBBsは\url{https://github.com/merlresearch/sebbs}の実装に変更を加えず，呼び出す形で使用しているため，省略する．
% その他，ソースコードは\url{https://github.com/custard-1855/dcase_base}にて公開している．\par

% また，使用したライブラリは以下の通り．
% \begin{itemize}
% \item 深層学習フレームワーク: PyTorch\cite{Ansel_PyTorch_2_Faster_2024}
% \item 音声前処理: torchaudio\cite{hwang2023torchaudio}
% \item モデル実装: PyTorch Lightning\cite{Falcon_PyTorch_Lightning_2019}
% \item ベースライン: DCASE 2024 Task 4\footnote{\url{https://github.com/DCASE-REPO/DESED_task/tree/master/recipes/dcase2024_task4_baseline}}
% \item 事前学習済みモデル: BEATs\cite{chen2022beats}
% \item 実験管理: Weights \& Biases (wandb)\footnote{\url{https://github.com/wandb/wandb}}
% \item 後処理: SEBBs\cite{ebbers2024sound}\footnote{\url{https://github.com/merlresearch/sebbs}}
% \item 評価指標計算: sed\_scores\_eval\footnote{\url{https://github.com/fgnt/sed_scores_eval}}，
% \item 数値計算: NumPy\cite{2020NumPy-Array}， pandas\cite{The_pandas_development_team_pandas-dev_pandas_Pandas}
% \end{itemize}
